% !TEX program = pdflatex
\documentclass[11pt,a4paper]{article}
\usepackage[margin=25mm]{geometry}
\usepackage[T1]{fontenc}
\usepackage[utf8]{inputenc}
\usepackage{lmodern}
\usepackage{amsmath,amssymb,amsthm,mathtools}
\usepackage{booktabs,array,tabularx}
\usepackage{microtype}
\usepackage[hidelinks,unicode]{hyperref}
\hypersetup{pdftitle={Continuous Identities, Formalization and Explicit Budgets for Binary Goldbach -- Synthesis V3.1},pdfauthor={Durand Serge},pdfsubject={Research synthesis with completed native numerical addendum},pdfkeywords={binary Goldbach, von Mangoldt, Mellin, projection, Lean 4}}
\setlength{\emergencystretch}{3em}
\setlength{\parskip}{3pt}
\setlength{\tabcolsep}{4pt}
\renewcommand{\arraystretch}{1.14}
\newcolumntype{Y}{>{\raggedright\arraybackslash}X}
\newtheorem{theorem}{Theorem}[section]
\newtheorem{proposition}[theorem]{Proposition}
\newtheorem{definition}[theorem]{Definition}
\theoremstyle{remark}
\newtheorem{remark}[theorem]{Remark}
\newcommand{\Lam}{\Lambda}
\newcommand{\Sing}{\mathfrak S}
\newcommand{\C}{\mathbb C}
\newcommand{\R}{\mathbb R}
\newcommand{\N}{\mathbb N}
\newcommand{\Z}{\mathbb Z}
\newcommand{\paid}{\textbf{Compiled Lean proof.}\space}
\newcommand{\source}{\textbf{SOURCE: not compiled.}\space}
\newcommand{\paper}{\textbf{PAPER: written argument.}\space}
\DeclareMathOperator{\Arg}{Arg}
\DeclareMathOperator{\Log}{Log}
\DeclareMathOperator{\sech}{sech}
\title{Continuous Identities, Formalization and Explicit Budgets\\
for Binary Goldbach\\[5pt]\large Synthesis V3.1 with Numerical Addendum}
\author{Durand Serge\\\small Independent researcher\\\small\texttt{seisner1967@gmail.com}}
\date{4 October 2026}
\begin{document}
\maketitle
\begin{abstract}
This synthesis continues the V2 monograph while preserving its balance
framework and notation. Research in round 22 transfers the additive coupling
to continuous traces, Mellin integrals and character projections.
The inventory frozen after independent batch 35 contains 88 modules and
1\,488 auxiliary declarations, including definitions. The new established
results include exact extraction of the coefficient of the genuine von
Mangoldt function on the circle, its alias-free discrete projection, canonical
logarithm quantization at scale $2^{58}$, the five finite-field roots specified
by the contract, and a Mellin representation with a justified interchange
of sum and integral. The weighted truncation envelopes, their uniform geometry
and their complete connection to the coefficient remain uncompiled proof
sources. An exact rational test validates a scalar upper bound below $10^{-6}$
at $N=10^8$; it does not evaluate the coefficient. A subsequent complete
native computation at $N=10^8$, $K=2^{27}$ reconstructs the canonical integer
$I_N$ by five-modulus CRT and confirms equality with the direct $A32$ sum
and the auxiliary $B40$ reconstruction. Their observed difference is zero,
within the fixed quantization budget $2E_{\log}<10^{-6}$. These observations
do not establish native program refinement in Lean or validate the Mellin
truncation chain. The formal inventory is unchanged in this numerical revision.
The target bound for $D_N$, the global signed contribution and the binary
Goldbach conjecture remain open.
\end{abstract}
\noindent\textbf{Keywords:} binary Goldbach; von Mangoldt; Mellin;
continuous projection; explicit errors; Lean 4.
\par\noindent\textbf{Classification:} 11P32.
\par\noindent\textbf{Scope.} This text closes a research stage and provides
a continuation contract. It announces neither a proof of Goldbach nor a
breakthrough of the parity barrier. The statuses refer to preserved receipts,
not to the number of identities that might still be formalized.
\par\noindent{\small The manuscript was prepared with assistance from research
agents and writing tools. This assistance does not constitute external
scientific validation. Compiled proofs are identified by their sources and
independent audits of their axioms.}

\clearpage
\section{Preserved framework and verification boundary}\label{sec:frame}
For an even integer $N$, we retain the notation
\[
u=\log N,\qquad L=1+u,\qquad \ell=\log u,
\]
with natural logarithms. The onset of the fixed analytic framework is
\begin{equation}\label{eq:onset}
u\ge 10^{24}.
\end{equation}
The point $N=10^8$ serves as a finite testing guide; it lies outside this regime.
A check at this point covers neither the onset \eqref{eq:onset} nor all even
integers below it. This stage preserves the established framework and examines
its connection to new continuous objects.

The reference quantity in the balance remains
\begin{equation}\label{eq:reference}
R_{\mathrm{ref}}=\Sing(N)(N-F_N),\quad
\Sing(N)=2C_2\prod_{\substack{p\mid N\\p>2}}\frac{p-1}{p-2},\quad
C_2=\prod_{p>2}\left(1-\frac1{(p-1)^2}\right).
\end{equation}
The earlier analytic dependencies of $F_N$ and the reference retain their
status in the monograph. This synthesis does not turn a reported analytic
input into a new self-contained Lean theorem. The research target remains
\begin{equation}\label{eq:target}
D_N\le \tau_N:=\frac{N}{256u\ell}.
\end{equation}
The cutoffs of the preserved profile are
\[
\alpha=\lceil N^{1/4}\rceil,\quad Q=\left\lfloor\frac{N-1}{\alpha}\right\rfloor,
\quad a_{\mathrm{cut}}=\lceil N^{7/16}\rceil,
\quad M_{\mathrm{cut}}=\lceil N^{3/4}\rceil.
\]
The thermal parameter $a>0$ introduced below is distinct from
$a_{\mathrm{cut}}$. All boundary faces, bands, corners, units and support
restrictions in the balance retain their positions.

\subsection*{Four levels of evidence}
\begin{tabularx}{\textwidth}{@{}p{29mm}Y@{}}
\toprule
Level & Exact meaning in this synthesis\\
\midrule
Compiled Lean & Frozen source, successful independent compilation, exact
catalogue and axiom audit; credit applies to the whole module.\\
SOURCE & Written and possibly reviewed Lean proof text; elaboration and
compilation are still absent or have failed.\\
PAPER & Identity or estimate derived in a note; its formalization and
evaluation obligations remain explicit.\\
Observed computation & Program actually run under a precise contract;
its result receives credit only within the scope of that contract.\\
\bottomrule
\end{tabularx}
The 1\,488 counted declarations include definitions. They are not 1\,488
Goldbach theorems. The standard axioms appearing in the audits include
\texttt{propext}, \texttt{Classical.choice} and \texttt{Quot.sound}.
No \texttt{sorryAx} is accepted in a credited module. Recovery output after
failed elaboration remains evidence of a technical failure, with no proof credit.

\clearpage
\section{The genuine thermal trace and the additive coefficient}\label{sec:heat}
We use the real arithmetic function
\[
\Lam(0)=\Lam(1)=0,\qquad
\Lam(n)=\begin{cases}\log p,&n=p^k,\ p\text{ prime},\ k\ge1,\\0,&\text{otherwise}.
\end{cases}
\]
All prime powers are retained. For $a>0$ and $\theta\in\R$,
\begin{equation}\label{eq:heat}
T_a(\theta)=\sum_{n\ge0}\Lam(n)e^{-an}e^{in\theta},\qquad
C_N=\sum_{n=0}^{N}\Lam(n)\Lam(N-n).
\end{equation}
No abstract trace function is needed: the series above is the object
actually formalized.

Set $r=e^{-a}\in(0,1)$ and
\begin{align}
U(a)&=\frac{r}{(1-r)^2},\label{eq:U}\\
B(a,M)&=r^{M+1}\left(\frac{M+1}{1-r}+U(a)\right).
\label{eq:heat-tail}
\end{align}
\paid Batch 11 directly constructs $0\le\Lam(n)\le n$, summability
of the terms, continuity of $T_a$, and the uniform bounds
\[
|T_a(\theta)|\le U(a),\qquad
|T_a(\theta)-T_{a,M}(\theta)|\le B(a,M),
\quad T_{a,M}(\theta)=\sum_{n=0}^{M}\Lam(n)r^ne^{in\theta}.
\]
It also proves continuity of the envelopes for $a>0$ and bounds the difference
between the full and partial integral projections by
$2e^{aN}U(a)B(a,M)$. The geometric majorant is constructed; it is not
supplied as a precision hypothesis.

\begin{theorem}[Exact continuous projection; batch 13]\label{thm:circle}
For every $a>0$ and every $N\in\N$,
\begin{equation}\label{eq:circle}
C_N=\frac{e^{aN}}{2\pi}\int_0^{2\pi}
T_a(\theta)\,\overline{T_a(-\theta)}\,e^{-iN\theta}\,d\theta
=\frac{e^{aN}}{2\pi}\int_0^{2\pi}T_a(\theta)^2e^{-iN\theta}\,d\theta.
\end{equation}
\end{theorem}
\paid The proof first extracts the characters on a finite rectangle,
establishes integral orthogonality and the normalization factor, then passes
to the limit using the envelope \eqref{eq:heat-tail}. Once $M\ge N$, the
integral projection of the polynomial $T_{a,M}$ already equals $C_N$ exactly.

Conjugation applies to the reflected phase:
$\overline{T_a(-\theta)}=T_a(\theta)$. The product $|T_a(\theta)|^2$
would select frequency differences; it does not yield the required additive
coefficient. The form centered on $[-\pi,\pi]$ relies on periodicity of the
genuine series. Its combination with the Mellin truncation is written
separately, with the SOURCE status described in Section~\ref{sec:radius}.

\clearpage
\section{Alias-free finite projector and finite-field projection}\label{sec:finite}
The continuous identity has a fully finite counterpart. It provides an exact
contract for the completed finite computation without replacing the genuine function
$\Lam$. For $K\ge1$, set $\theta_j=2\pi j/K$.

\begin{theorem}[Discrete orthogonality and extraction; batch 15]
For $N\le M$ and
\begin{equation}\label{eq:alias}
K>\max(N,2M-N),
\end{equation}
for every real $a$, we have
\begin{equation}\label{eq:discrete}
C_N=\frac{e^{aN}}{K}\sum_{j=0}^{K-1}
T_{a,M}(\theta_j)^2e^{-iN\theta_j}.
\end{equation}
\end{theorem}
\paid Batch 15 proves orthogonality of the actual exponential characters,
then their expansion and selection of the antidiagonal. Condition
\eqref{eq:alias} excludes spurious frequencies $m+n-N\in K\Z$ with
$0\le m,n\le M$. It is a mathematical guard, not an experimental tolerance.
The choice $a=0$ is valid for this finite polynomial; no infinite trace
$T_0$ is introduced.

An analogous projection exists in a finite field. For a prime $p$, an element
$\omega\in\mathbb F_p$ of exact order $K$ and $K\ne0$ in $\mathbb F_p$,
orthogonality gives
\[
\frac1K\sum_{j=0}^{K-1}\omega^{jk}
=\begin{cases}1,&K\mid k,\\0,&K\nmid k.\end{cases}
\]
\paid Batch 19 proves this identity from the primitive root;
the resulting orthogonality is not a premise. For a finite sequence $A_n$
and \eqref{eq:alias}, its projection selects
$\sum_{n=0}^N A_nA_{N-n}$ in $\mathbb F_p$.

The continuous-circle representation and this character projection share
one object: the exact additive coefficient. A projection identity does not
show that this coefficient is sufficiently large. Nor does it establish
that a fast-transform program implements the theorem's finite sum.
That correspondence requires stage invariants, normalizations, machine-word
bounds and integer reconstruction.

The final contract fixes
\begin{equation}\label{eq:fixed}
N=M=100\,000\,000,\qquad K=2^{27}=134\,217\,728,
\qquad S=2^{58}=288\,230\,376\,151\,711\,744.
\end{equation}
Here $K>N$ satisfies \eqref{eq:alias}. The finite proofs remain valid
regardless of whether a computation of this size is completed. The validation
contract must retain coefficients from all prime powers, including those
of exponent at least two.

\clearpage
\section{Canonical logarithms and explicitly constructed precision}\label{sec:quant}
The quantization in batch 17 does not assume oracle-supplied precision.
For $2\le p\le10^8$, set
\[
k=\lfloor\log_2p\rfloor,\qquad
z=\frac{p-2^k}{p+2^k},\qquad 0\le z<\frac13,
\]
and define the rational sum
\begin{equation}\label{eq:log32}
L_{32}(z)=2\sum_{j=0}^{31}\frac{z^{2j+1}}{2j+1},\qquad
\rho_{32}=\frac{9}{4\cdot65\cdot3^{65}}.
\end{equation}
The real series representation of $\log(1+z)-\log(1-z)$, its geometric
tail and the decomposition using $2^k$ give the interval
\begin{equation}\label{eq:logbox}
b_p=kL_{32}(1/3)+L_{32}(z),\qquad
\log p\in[b_p,b_p+(k+1)\rho_{32}].
\end{equation}
\paid On the fixed domain, the width satisfies
\[
(k+1)\rho_{32}\le\frac{243}{260\cdot3^{65}}<2^{-96}.
\]
The canonical point $A_p$ is obtained by multiplying the midpoint of
\eqref{eq:logbox} by $S$, rounding to the nearest integer with ties to even,
then clamping to $[0,32S]$. The constructor uses exact rational quotient
and remainder. It is not identified with a different rounding function
without proof.

For all integers, define
\[
A_n=\begin{cases}
A_{\min\mathrm{Fac}(n)},& n\text{ is a prime power},\\
0,&\text{otherwise},
\end{cases}\qquad q_n=A_n/S.
\]
The cases $0$ and $1$ are handled separately. Using the direct definition
of $\Lam$ and its least prime factor retains every $p^k$.
\paid The proof establishes, for $0\le n\le10^8$,
\begin{equation}\label{eq:q-error}
0\le\Lam(n),q_n\le32,\qquad |\Lam(n)-q_n|\le 1/S.
\end{equation}
Precision follows from \eqref{eq:log32}--\eqref{eq:logbox} and the constructed
rounding rule; it is not a free hypothesis of the theorem.

Set $I_N=\sum_{n=0}^{N}A_nA_{N-n}\in\Z_{\ge0}$. The same batch proves
\begin{equation}\label{eq:E-log}
\left|C_N-\frac{I_N}{S^2}\right|\le E_{\log}(N,S),\qquad
E_{\log}(N,s)=(N+1)\left(\frac{64}{s}+\frac1{s^2}\right),\quad s>0.
\end{equation}
The envelope is continuous for $s>0$. The fixed integer guard
\begin{equation}\label{eq:q-guard}
2(N+1)(64S+1)10^6\le S^2
\end{equation}
is compiled, hence $2E_{\log}(10^8,2^{58})\le10^{-6}$. This result proves
the real quantization bound without computing $I_N$ or $C_N$.

\clearpage
\section{Five prime moduli, exact residues and native-code obligations}\label{sec:crt}
The contract fixes the following five pairs; their primality and the root
orders were independently compiled in batch 22.
\begin{center}
\begin{tabular}{@{}rrrr@{}}
\toprule
$i$ & $p_i$ & $(p_i-1)/K$ & $g_i$\\
\midrule
1 & 2\,013\,265\,921 & 15 & 31\\
2 & 2\,281\,701\,377 & 17 & 3\\
3 & 3\,221\,225\,473 & 24 & 5\\
4 & 3\,489\,660\,929 & 26 & 3\\
5 & 3\,892\,314\,113 & 29 & 3\\
\bottomrule
\end{tabular}
\end{center}
In $\mathbb F_{p_i}$, set $\omega_i=g_i^{(p_i-1)/K}$. The proof establishes
order $K$ from primality certificates and the required power checks.

\begin{theorem}[Concrete canonical projection; batch 22]
With parameters \eqref{eq:fixed} and the canonical integers $A_n$,
\begin{equation}\label{eq:ntt-projection}
F_i(j)=\sum_{n=0}^{N}A_n\omega_i^{jn},\qquad
r_i=K^{-1}\sum_{j=0}^{K-1}F_i(j)^2\omega_i^{-Nj}
=I_N\pmod {p_i}.
\end{equation}
\end{theorem}
\paid The theorem instantiates the finite-field projection with the
actual roots in the table. Neither the final residue nor coefficient
precision is taken as a premise.

\paper For $\Pi=\prod_{i=1}^5p_i$, the contract bounds are
\begin{equation}\label{eq:crt-size}
0\le I_N\le(N+1)(32S)^2<2^{153}<\Pi.
\end{equation}
They provide the domain for unique reconstruction by the Chinese remainder
theorem. A reconstruction $c\in[0,\Pi)$ with the five correct residues must
therefore equal $I_N$, provided the exact connection to the constructor and
projections is established. Proving the behavior of the CRT code, its machine
words and the fast transform remains separate from \eqref{eq:ntt-projection}.

The native producer and checker were both built in BUILD04: two compilations
completed with exit code zero. This establishes the existence of two frozen
executable images without proving their refinement of the Lean functions.
The completed operational checker verifies the entire catalogue up to $N$,
the classification of prime powers, equality of every recorded point to the
canonical $A32$ constructor, residues, reconstruction and integer guards
(Section~\ref{sec:numeric}). These are observed native checks, not a Lean
program-refinement theorem.
An auxiliary check using a $B40$ series is not the primary $A32$ constructor;
its code and real precision have no complete Lean proof at this stage.

The resource window is an operational constraint. Declared trust in Windows,
GCC and the native images is neither universal loader closure nor observation
of all dynamic dependencies. None of these operational assumptions replaces
a program-refinement theorem.

\clearpage
\section{Mellin: exact identity with the genuine von Mangoldt series}\label{sec:mellin}
Let $w\in\C$ satisfy $\Re w>0$. The principal logarithm determines
$w^{-s}=\exp(-s\Log w)$. Set
\begin{equation}\label{eq:mellin-objects}
s(t)=2+it,\quad
D(t)=\sum_{n\ge2}\Lam(n)n^{-2-it},\quad
P(w)=\sum_{n\ge0}\Lam(n)e^{-nw},\quad
K(w,t)=\Gamma(2+it)w^{-2-it}.
\end{equation}
The index-zero and index-one terms of $D$ vanish in the formal definition.
For $n\ge1$, the powers use the real logarithm of the positive number $n$.
In particular, $P(a-i\theta)=T_a(\theta)$.

\begin{theorem}[Inversion and infinite interchange; batch 35]\label{thm:mellin}
For $\Re w>0$, the function $t\mapsto D(t)K(w,t)$ is integrable and
\begin{equation}\label{eq:mellin}
P(w)=\frac1{2\pi}\int_{\R}D(t)\Gamma(2+it)w^{-2-it}\,dt.
\end{equation}
Moreover, $D$ is continuous and $|D(t)|\le6$ for every $t\in\R$.
\end{theorem}
\paid Batch 35 contains 30 declarations: 24 theorems and 6 definitions.
Its independent compilation and exact axiom audit succeeded. The interchange
does not rely on an assumed final integrability statement or an externally
supplied equality.

The proof begins with the positive series
\[
b_n=\Lam(n)n^{-2},\qquad
0\le b_n\le 2n^{-3/2}\ (n\ge1),\qquad \sum_n b_n\le6.
\]
It constructs uniform convergence of the series $D$ and its bound.
The Gamma inversion identities, holomorphy on the right half-plane and
integrable kernel majorants come from independently compiled dependencies.
For each $n>0$, the branch identity $(nw)^{-s}=n^{-s}w^{-s}$ is established
with $n$ a positive real number. Finally, summability of the integrals of
norms justifies the actual interchange between the series in $n$ and the
integral in $t$. The terms $\Lam(p^k)$ remain present.

On the line $\Re s=2$, the classical identification of this series with
$-\zeta'(s)/\zeta(s)$ belongs to a separate zeta connection. Theorem
\ref{thm:mellin} directly uses the genuine series $D$; its credit does not
validate all explicit formulas in zeros, all contour shifts or a uniform
evaluator for $\zeta$.

Identity \eqref{eq:mellin} transfers the coupling exactly into a continuous
representation. It does not yet give a favorable sign for coefficient
\eqref{eq:circle}. Absolute domination ensures existence of the integrals
and the interchanges; a new signed cancellation remains a separate obligation.

\clearpage
\section{Signed truncations and right-half-plane geometry}\label{sec:tails}
We retain the concrete definitions of the truncation packages. For
$\Re w>0$, let
\begin{equation}\label{eq:Cdelta}
\eta(w)=\frac{\pi/2+|\Arg w|}{2},\quad
\delta(w)=\frac{\pi/2-|\Arg w|}{2}>0,\quad
C(w)=|w|^{-2}\sec^2\eta(w).
\end{equation}
The Gamma kernel has constructed exponential majorants on this domain.
Independent, compiled batch 30 establishes the unweighted tails in both
orientations. The bridge in batch 32 connects full Gamma inversion to the
corresponding truncation. These credits do not automatically include the
new tail weighted by $D$.

\source The $E_\Lam$ package defines, for $H\ge0$,
\begin{align}
P_H(w)&=\frac1{2\pi}\int_{-H}^{H}D(t)K(w,t)\,dt,\\
T_+(w,H)&=\int_H^\infty D(t)K(w,t)\,dt,\qquad
T_-(w,H)=\int_H^\infty D(-t)K(w,-t)\,dt.
\end{align}
It aims to establish the exact statements
\begin{equation}\label{eq:lambda-tail}
P(w)-P_H(w)=\frac{T_-(w,H)+T_+(w,H)}{2\pi},\qquad
|P(w)-P_H(w)|\le
\frac{6C(w)e^{-\delta(w)H}}{\pi\delta(w)}.
\end{equation}
The factor 6 comes from Theorem \ref{thm:mellin}. Reflection applies to both
factors $D$ and $K$. The source contains concrete integrability and tail
proofs; no premise equal to \eqref{eq:lambda-tail} is added.

Batch 36 fails on a reflection rewrite at line 118:6. It produced no compiled
proof object for this module; the following geometry module was not invoked.
The repair is preserved in a new SOURCE02 package, with no further compilation
at the close of this stage. The diagnostic concerns elaboration, without
an analytic refutation or a proof of a parity obstruction.

\source The Geometry25 package constructs, for $a>0$,
\begin{equation}\label{eq:geometry}
C(a-i\theta)=\frac2{a^2}
\left(1+\frac{|\theta|}{\sqrt{a^2+\theta^2}}\right)\le\frac4{a^2}.
\end{equation}
For $|\theta|\le\pi$, it also aims to establish
\begin{equation}\label{eq:d}
\delta(a-i\theta)\ge d(a):=\tfrac12\arctan(a/\pi)>0.
\end{equation}
Branches, quadrants, signs, square roots and continuity are explicitly
written out. The radii are closed expressions and continuous on their domains:
$\Re w>0$ and $H\in\R$ for the expression, $H\ge0$ for the tail bound.
Continuity of an expression is not an additional proof of continuity of
an integral with a moving boundary.

\clearpage
\section{Coefficient envelope and remaining scalar connection}\label{sec:radius}
From \eqref{eq:lambda-tail}--\eqref{eq:d}, the proposed uniform radius is
\begin{equation}\label{eq:eps}
\varepsilon(a,H)=\frac{24e^{-d(a)H}}{\pi a^2d(a)},\qquad a>0,\ H\ge0.
\end{equation}
The coefficient source defines the centered projection of the truncation
\[
C_{N,H}=\frac{e^{aN}}{2\pi}\int_{-\pi}^{\pi}
P_H(a-i\theta)^2e^{-iN\theta}\,d\theta
\]
and the closed-form radius
\begin{equation}\label{eq:E-mellin}
\mathcal E_N(a,H)=e^{aN}\bigl(2U(a)\varepsilon(a,H)
+\varepsilon(a,H)^2\bigr).
\end{equation}
\source The 34 declarations in this package construct transport to the
centered circle, angular continuity and integrability at fixed $H$, then
aim to establish
\begin{equation}\label{eq:coef-error}
|C_N-C_{N,H}|\le\mathcal E_N(a,H).
\end{equation}
The elementary inequality used is
$|x^2-y^2|\le2|x||x-y|+|x-y|^2$. Together with $|T_a|\le U(a)$, it explains
the two budget terms. Joint continuity of the radius is written for $a>0$
and $H\in\R$; its interpretation as a truncation error requires $H\ge0$.
The theorem is not compiled at stage closure.

The genuine series $T_a$ is periodic. The complex power
$(a-i\theta)^{-2-it}$, taken alone on its principal branch, is not.
The source does not postulate periodicity of $P_H$. Shifting the projection
interval concerns the full trace, whose periodicity follows from the series;
the truncation then receives its own budget.

\source The scalar package of 26 declarations defines the same functions
$U$, $d$, $\varepsilon$ and $\mathcal E_N$, and writes a real-variable
chain at the point
\begin{equation}\label{eq:point}
N=10^8,\qquad a=10^{-8}=1/N,\qquad H=10^{11},\qquad\tau=10^{-6}.
\end{equation}
It aims in particular to prove $d(a)>1/(8N)$, a bound for $U$ through the
$\sinh$ identity, $e^{125}>10^{50}$, then
$\mathcal E_N(a,H)\le576/10^{10}<\tau$.
It is not compiled. A separate bridge of eight theorems writes the equalities
between the two radius definitions and the instantiation of
\eqref{eq:coef-error} by this scalar result. Its two imports are prospective;
the bridge is also uncompiled.

These sources address precise theoretical obligations without establishing
a numerical coefficient. Before their independent compilation, no complete
Lean chain supports promoting \eqref{eq:coef-error} and its instantiation
at \eqref{eq:point} to an established formal result.

\clearpage
\section{The completed rational test and its scope}\label{sec:scalar-test}
The observed scalar test uses only exact fractions. At \eqref{eq:point},
it defines
\[
S_{150}(125)=\sum_{k=0}^{150}\frac{125^k}{k!},\qquad
q=\frac1{S_{150}(125)},\qquad
\mathcal B=3\bigl(128N^5q+4096N^6q^2\bigr).
\]
\textbf{Observed computation: rational PASS.} One actual invocation,
with no retry, verifies exactly
\begin{equation}\label{eq:B-test}
\mathcal B<10^{-6}.
\end{equation}
The preserved checks include
$S_6(5)>100$ and
$S_{150}(125)\ge S_6(5)^{25}>10^{50}$.
The result is linked to the executor's receipt and the ROOT physical
observation listed in the appendix. It uses no floating-point approximation.

Comparison \eqref{eq:B-test} is a scalar result. The program does not integrate
$D$, compute $P_H$ or $C_{N,H}$, or evaluate $C_N$. It does not verify a prime
catalogue, the native transform or a Goldbach identity. The connection between
$\mathcal B$ and the real radius \eqref{eq:E-mellin} uses the analytic
inequalities written in the scalar package and the bridge of Section
\ref{sec:radius}. Their status remains SOURCE.

\paper The analytic connection to the same radius is explicitly as follows.
The real inequalities in the connectors give $3<\pi<7/2$,
$d(1/N)>1/(8N)$ and
$U(1/N)=1/(4\sinh^2(1/(2N)))\le N^2$.
Thus $dH>125$ at the fixed point. The positive Taylor series gives
$S_{150}(125)\le e^{125}$, hence $e^{-dH}\le e^{-125}\le q$.
Expression \eqref{eq:eps} therefore satisfies
\[
\varepsilon(1/N,H)\le
\frac{24}{3}N^2(8N)q=64N^3q.
\]
Since $aN=1$ and $e<3$, \eqref{eq:E-mellin} implies
\[
\mathcal E_N(1/N,H)\le
3\bigl(2N^2\cdot64N^3q+(64N^3q)^2\bigr)=\mathcal B<\tau.
\]
This chain explains the proposed analytic scope of the test exactly.
It remains PAPER/SOURCE until compilation of the connectors and the bridge
to the coefficient; the rational computation alone does not verify
these dependencies.

\subsection*{Feasibility of a literal Mellin evaluation}
\paper The sharper geometry gives the uniform radius
\[
\overline C(a)=\frac2{a^2}\left(1+
\frac{\pi}{\sqrt{a^2+\pi^2}}\right)<\frac4{a^2},\qquad
\overline\varepsilon=
\frac{6\overline C(a)e^{-d(a)H}}{\pi d(a)}.
\]
The cost note studies the threshold of the majorant
$e^{aN}(2U\overline\varepsilon+\overline\varepsilon^2)$.
For parameters \eqref{eq:point}, it places this threshold between
$5\cdot10^{10}$ and $10^{11}$ using closed-form inequalities.
This estimate concerns a sufficient condition obtained from a majorant;
it is not a lower bound for the actual error and does not impose this
height on every possible method.

The same note bounds a literal tail of $D$ on paper by $(\log M+1)/M$
and angular quadrature through a derivative majorant. Extremely crude
choices such as $M=10^{52}$ and $J=10^{67}$ illustrate symbolic closure
of these budgets. They describe neither an available program nor a run
that could finish within this stage's resource window. Directed certification
of the primitives at every node and vertical quadrature remain open.
The scalar success \eqref{eq:B-test} is therefore not a performance
measurement of a global Mellin evaluation.

\clearpage
\section{Completed native evaluation and its verification boundary}\label{sec:numeric}
The finite contract in Section \ref{sec:crt} gives a falsifiable target:
produce $I_N$ from the canonical coefficients, reconstruct the same integer
from the five residues, and verify all guards. Then \eqref{eq:E-log}
provides the interval
\[
C_N\in\left[I_N/S^2-E_{\log},\ I_N/S^2+E_{\log}\right].
\]
The radius is not chosen after observation; its real construction and guard
are already compiled. The completed native checker verifies correspondence
of the records to $A32$, complete catalogue coverage and the integer
computation. Refinement of its machine implementation to the Lean objects
remains a separate obligation.

\paper For $K=2^{27}$, a radix-2 transform has
$K\cdot27/2=1\,811\,939\,328$ butterflies per transform. The contract counts
give tens of billions of modular products for the five moduli and their
independent check. An array of $K$ 32-bit words occupies 536\,870\,912 bytes;
the 64-bit points on $0,\ldots,N$ occupy 800\,000\,008 bytes, and the 32-bit
factors occupy 400\,000\,004 bytes. These payload counts are neither a total
memory measurement nor a runtime promise.

\textbf{Historical V3 observation: STOP without a verdict.} Consumer05 reused
three frozen producer04 payloads as identical copies and a checker image
from BUILD04. The payloads were not validated by that stopped attempt.
The attempt creates one checker, no producer and no compiler,
with no retry. The parent budget is 3\,600 seconds; the checker has a deadline
at 3\,300 seconds, leaving a 300-second reserve for logs and preservation checks.

The checker stops with code 238 and status \texttt{HARD\_WALL\_NO\_VERDICT}.
The parent exits with code 1 after approximately 3\,312 seconds, or 55 minutes.
No coefficient was present in that receipt. Closure documents quiescence and the POST actually
produced; the ROOT observation verifies preservation of 6\,576 inputs,
109 capture pairs, three payload pairs and 3\,089 archive files.
The stop is a resource failure without a mathematical verdict, not a
counterexample to the identity or the constructed precision.

Those OS resource controls distinguish a 2 GiB Job commit-memory limit from
RSS monitoring, which is not an OS-enforced RSS limit. Operational trust
in Windows/GCC/native code remains declared and bounded in scope; universal
loader closure is not asserted. No new attempt was included in the V3
closure. These historical receipts remain unchanged.

\textbf{New observed computation: exact numerical arithmetic validated.}
On 4 October 2026, a fresh execution of the unchanged BUILD04 producer and
the operational checker both finish with exit code zero. The checker retains
the reference $A32$ construction and uses an exact integer Horner evaluation
of its rational logarithm series, with a cached rational series for $\log 2$.
A 128-case comparison of the exact lower and upper bounds against the
reference implementation passes before the full run. The catalogue,
prime-power classification, DIF transforms, CRT and direct integer sum
remain within the same arithmetic contract.

The checker validates 5\,761\,455 prime records, all canonical words and the
five residues. CRT reconstruction, the direct $A32$ sum and auxiliary $B40$
reconstruction agree exactly. The fresh producer's three payload files are
byte-identical to the checked copies. Appendix~\ref{app:numerical} records
the integer, error radius, resource measurements and receipt bindings.
The independent Python audit reconstructs CRT and checks receipts and hashes;
it does not independently replay the full prime catalogue or transforms.

The operational Windows backend avoids privileged working-set quotas,
the source of error 1314, and requests a 3 GiB Job commit-memory limit.
RSS is monitored, not enforced as an OS working-set quota.
The new checker and parent have a 24-hour wall budget; the unchanged producer
finishes within its original internal deadline. An initial startup stop due
to the console helper exceeding a one-process monitor is preserved;
the completed invocations permit 16 active and 32 total Job processes.
Neither completed launcher invocation performs an automatic retry.

\clearpage
\section{Connection to the balance: the signed obligation remains}\label{sec:dn}
The coefficient $C_N$ includes proper prime powers.
Define directly
\[
P_{\mathrm{pr}}(n)=\begin{cases}\log n,&n\text{ prime},\\0,&\text{otherwise},\end{cases}
\qquad V_{\mathrm{pp}}(n)=\Lam(n)-P_{\mathrm{pr}}(n).
\]
The notation $V_{\mathrm{pp}}$ avoids confusing this weight with the amplitude
$U(a)$. In the ledger audit, the coefficient of prime pairs and the
correction are
\begin{align}
R_{\mathrm{sf}}&=\sum_{n=0}^N P_{\mathrm{pr}}(n)P_{\mathrm{pr}}(N-n),\\
Q_N&=C_N-R_{\mathrm{sf}}\nonumber\\
&=\sum_{n=0}^N\bigl(2P_{\mathrm{pr}}(n)V_{\mathrm{pp}}(N-n)
+V_{\mathrm{pp}}(n)V_{\mathrm{pp}}(N-n)\bigr)\ge0.
\label{eq:PP}
\end{align}
\paper The direct upper bound $Q_N\le2\sqrt N\,u^2$ for $N\ge2$
is retained in this note. It is not credited here as a new
Lean theorem. This global term is not interchangeable with the
prime-power terms whose supports are separately retained in the balance.

The monograph's bridge, within its stated scope, under its conditions and with exactly
the same support, retains
\begin{equation}\label{eq:ledger}
D_N+\varepsilon_{\mathrm{tr}}
=R_{\mathrm{ref}}-R_{\mathrm{sf}}+2\max(e,0).
\end{equation}
Its connection to the continuous coefficient is therefore
\begin{equation}\label{eq:gap}
D_N=R_{\mathrm{ref}}-C_N+Q_N+2\max(e,0)-\varepsilon_{\mathrm{tr}}.
\end{equation}
The ledger retains each of
$B_{\mathrm{prime}}^{(a_{\mathrm{cut}})}$,
$B_{\mathrm{pp}}^{(a_{\mathrm{cut}})}$,
$P_{\mathrm{band},\ge2}$, $Z_{\mathrm{face},\ge2}$, $I_\alpha$ and
$2\max(e,0)$ exactly once. This transposition does not redevelop these objects
using the arithmetic methods abandoned in loop 22.

Knowing a narrow interval for $C_N$ does not give the sign
of $R_{\mathrm{ref}}-C_N$, or control of $e$, the boundaries and the corners.
To establish \eqref{eq:target}, a signed estimate
compatible with \eqref{eq:gap} remains to be proved, retaining these obligations and the
bridge conditions. Assuming this estimate as a premise equivalent to
the target would not constitute a way around the parity barrier.

The continuous identities establish exact representations, and the
budgets describe a truncation or quantization error. No
global cancellation of the true spectral contribution transferred to the
ledger has been proved. This obligation remains open after the present
finite computation of $I_N$ at $10^8$ and the enclosure for $C_N$ under the
native correspondence assumptions. It must be addressed at the
threshold \eqref{eq:onset} and with the finite coverage required by the framework.

\clearpage
\section{Boundaries, zeros and invariance: directions with precise scope}\label{sec:strategy}
The continuous pivot preserves phase before taking a norm. For
$a>0$, $N\ge1$ and $q\in\C$, consider
\[
J_{a,N}(q)=\frac1{2\pi}\int_{-\pi}^{\pi}
(a-i\theta)^{-q}e^{-iN\theta}\,d\theta.
\]
\paid Batch 25 formalizes integration by parts and the true boundary term
on the principal branch:
\begin{equation}\label{eq:bord}
J_{a,N}(q)=\frac{i(-1)^N}{2\pi N}
\bigl((a-i\pi)^{-q}-(a+i\pi)^{-q}\bigr)
+\frac qN J_{a,N}(q+1).
\end{equation}
The boundary term does not vanish pointwise. This module of 30 declarations
provides no favorable sign for the full spectral projection.
Periodicity of the trace after inversion does not make
each truncated kernel periodic.

\paper The signed explicit-formula strategy localizes a smooth test
function near $x+y=N$, then aims for the decomposition
\[
C_N=\mathcal B_N-2\mathcal L_N+\mathcal Z_N.
\]
The three terms respectively denote the background, coupling to the
zero channel, and the two-zero contribution in the proposed explicit formula.
The full identity, archimedean corrections, multiplicities,
interchanges and zero tails remain to be established in this channel.
The retained reviews find consistency subject to these obligations; they
credit neither a sign for $\mathcal Z_N-2\mathcal L_N$ nor a
Lean proof of the global formula.

Organization into conjugation--reflection orbits retains zeros off
the critical line. A four-point orbit generates a density
of the form
\[
4m x^{-1/2}\cosh\bigl((\beta-1/2)\log x\bigr)
\cos(\gamma\log x),
\]
and a two-point critical orbit gives the density
$2m x^{-1/2}\cos(\gamma\log x)$, with multiplicity $m$.
Stabilizers prevent double counting a degenerate orbit.
The cosine phase can change sign: correct organization of the
orbits does not automatically give positivity of the additive coefficient.

The review of the canonical kernel exhibits a negative direction in a
two-dimensional restriction; it rules out the naive positivity argument
examined there. Its scope is this specific kernel. It is neither a
universal impossibility result for every operator approach nor a
refutation of Goldbach. Likewise, any proposed modularity or exact
invariance must be proved for the true objects and must
isolate the signed term required in \eqref{eq:gap}. No such
sufficient connection is presented as imminent at this closure.

\clearpage
\section{Continuous Fej\'er: reconstruction and the remaining sign}\label{sec:fejer}
The new spectral note starts from the true series $D$ of batch 35, with
$c_n=\Lam(n)/n^2$ for $n\ge2$ and $c_0=c_1=0$.
The mass bound $\sum c_n\le6$ makes $D$ positive definite: for finite
families $(z_j,t_j)$,
\[
\sum_{j,k}\overline z_jz_kD(t_k-t_j)
=\sum_{n\ge2}c_n\left|\sum_k z_ke^{-it_k\log n}\right|^2\ge0.
\]
\paper For $T>0$ and an integer $m\ge1$, the triangular average
\begin{equation}\label{eq:fejer-cont}
\begin{split}
A_T(m)&=\frac1T\int_{-T}^{T}\left(1-\frac{|t|}{T}\right)
D(t)e^{it\log m}\,dt\\
&=\sum_{n\ge2}c_n\operatorname{sinc}^2\!\left(\frac T2\log\frac mn\right)
\end{split}
\end{equation}
is real and nonnegative. Here $\operatorname{sinc}(x)=\sin(x)/x$ if
$x\ne0$ and $\operatorname{sinc}(0)=1$. The interchange over the compact interval is
justified by summability; the triangle is the autocorrelation of an
interval indicator. Distinct frequencies satisfy
$|\log(m/n)|\ge\log(1+1/m)$. The note therefore derives the closed radius
\begin{equation}\label{eq:fejer-radius}
0\le A_T(m)-c_m\le
\frac{24}{T^2\log^2(1+1/m)},\qquad T>0.
\end{equation}
For $m=1$, $c_1=0$ and the same bound holds. The radius is continuous
for $T>0$ and tends to zero for fixed $m$. Neither this identity nor its
connecting results are compiled in this stage.

This mechanism reconstructs amplitudes from the true vertical data,
then can insert them into the reflected blocks $n,N-n$. It does not create
shared mass for a pair. The note retains a bounded countermodel:
the sector consisting solely of powers of 2 is positive definite and satisfies the
same reconstruction, while its additive coefficient at $N=14$ is zero.
This observation rules out deducing a uniform strict lower bound from
this positivity alone. It does not replace the true zeta function and does not rule out an
additional identity specific to its full data.

A future evaluator would have to produce intervals for the true
integrals $A_T(m)$, with controlled primitives, locations and quadratures.
\eqref{eq:fejer-radius} does not account for any of these evaluation
errors. The cost increases with $m$, and no feasibility at $10^8$
is asserted. The earlier identities for divisible grids and their
covariances remain in V2, with their support conventions; this
direction concerns the continuous spectral variable $t$.

Research is frozen. Any resumption must establish an additional signed
connection to the ledger \eqref{eq:gap}, rather than rename this
reconstruction as parity cancellation. No mechanism sufficient
for \eqref{eq:target} or imminent discovery is claimed.

\clearpage
\section{Frozen formal inventory and record of limitations}\label{sec:inventory}
The independent observation of batch 35 fixes the official inventory at
\textbf{88 modules / 1\,488 auxiliary declarations, including definitions}.
Batch 36 adds no credit. The following table isolates the elements
used in this synthesis; it is not an exhaustive total of the
historical inventory.

\begingroup\small\setlength{\tabcolsep}{3pt}
\noindent
\begin{tabularx}{\textwidth}{@{}p{12mm}Yp{22mm}p{25mm}@{}}
\toprule
Batch & Object & Declarations & Status\\
\midrule
11 & Geometric envelope of the true thermal trace & 42 & Row PASS ; batch FAILED\\
13 & Continuous extraction identity & 20 & Lean PASS\\
15 & Orthogonality and discrete complex projector & 29 & Lean PASS\\
17 & Canonical logarithms and envelope & 52 & Lean PASS\\
19 & Generic finite-field projector & 24 & Lean PASS\\
22 & Five concrete roots and $A32$ projection & 28 & Lean PASS\\
25 & Integration by parts and boundary term & 30 & Lean PASS\\
26/27 & Local and holomorphic Gamma inversion & 22/15 & Row PASS26 ; PASS27\\
30/32 & Gamma tails and unweighted bridge & 22/2 & Lean PASS\\
35 & Mellin transform of the true $\Lam$ series & 30 & Lean PASS\\
36 & Weighted tail $E_\Lam$ ; geometry & 16 ; 25 & FAIL ; not invoked\\
\bottomrule
\end{tabularx}
\endgroup

Batches 11 and 26 have a global FAILED status: each retains one
whole-module PASS row (Envelope42 and Local22, respectively),
while another module in the same batch fails. Their partial credits
are attested by the adjudications and observations; the failed
modules are not counted as established. Batches 13 and 27
later validate distinct corrected sources.

The coefficient sources (34), scalar connecting results (26) and the
coefficient--scalar bridge (8) have not been compiled. Their declaration
counts are source catalogues; they do not add to
1\,488. Earlier \texttt{SOURCE ONLY} comments in some
files that subsequently became PASS remain in the historical captures; the current
status comes from the adjudication and receipt, not from retroactive rewriting.

Correction batches have preserved their failures: batch 16 contains
series and rational-cast connecting issues, repaired and compiled
in batch 17; batches 31 and 33 contain instance
or continuity issues, repaired in the source actually compiled in batch
35. Batch 36 fails on reflection of the product of two functions.
These diagnostics are not proofs of the parity barrier. Conversely,
an auxiliary PASS is not a victory over that obstruction.

Verification is a property of the exact file, exact catalogue
and its dependencies. The checks retain earlier logs,
captures, sources and archives; they distinguish full reading of a
text, targeted reading, metadata inventory and byte hashing.
The historical companion file \path{evidence_v3.json} specifies the V3
boundary. The new \path{evidence_v3_1.json} binds this revision to its
completed numerical receipts without modifying the historical registry.

\clearpage
\appendix
\section{Compact provenance and resumption}\label{app:provenance}
All paths below are relative to the corpus root
\path{Goldbach_Parity_20261002}. SHA-256 prefixes aid reading;
complete hashes, sizes, statuses and reading scopes are
retained in \path{synthesis_v3/evidence_v3.json}. Source paths
do not replace execution receipts.
The published corpus is available in the repository
\href{https://github.com/seisner1967-hash/goldbach-horizon}{\texttt{goldbach-horizon}},
with the V3 companion registry.

\begingroup\small\setlength{\tabcolsep}{3pt}
\noindent
\begin{tabularx}{\textwidth}{@{}p{20mm}Yp{23mm}@{}}
\toprule
Item & Relative path & SHA prefix\\
\midrule
Mellin PASS & \path{round22/judge5/batch35/sources/ComplexGammaMellinLambda22.lean} & \texttt{4a94415678e4}\\
Receipt 35 & \path{round22/judge5/batch35/batch35_attempt01/receipt.json} & \texttt{3cffa331b392}\\
Observation 35 & \path{.arbor/sessions/parity/.coordinator/messages/round22_judge_batch35_closed_observation.json} & \texttt{dc5985976051}\\
Repaired tail & \path{round22/role4/complex_gamma_mellin_lambda_tail_source02/ComplexGammaMellinLambdaTail22.lean} & See JSON\\
Geometry SOURCE & \path{round22/role4/complex_gamma_circle_geometry_source01/ComplexGammaCircleGeometry22.lean} & \texttt{2322915060c6}\\
Coefficient SOURCE & \path{round22/role4/lambda_circle_truncation_source01/LambdaCircleTruncationEnvelope22.lean} & \texttt{bf9b8257760a}\\
Scalar SOURCE & \path{round22/role4/scalar_radius_real_connectors_source01/ScalarRadiusRealConnectors22.lean} & \texttt{6ba4b8d4cca0}\\
Bridge SOURCE & \path{round22/role3/lambda_circle_scalar_bridge_source01/LambdaCircleScalarRadiusBridge22.lean} & \texttt{cfb0d2cfd238}\\
Rational point & \path{round22/role6/mellin_circle_radius_point_source01/actual_radius_point01_attempt01/receipt.json} & \texttt{d46e3ffeb6d0}\\
Ledger & \path{round22/role4/continuous_dn_gap_paper01/dn_gap_audit22.md} & See JSON\\
\bottomrule
\end{tabularx}
\endgroup

\subsection*{Self-contained handoff contract}
The historical document \path{synthesis_v3/HANDOFF_THEORY_V3.txt} is preserved.
Its successor \path{HANDOFF_FEJER_V3_1.txt} records the completed finite
integer computation, receipt locations and unchanged proof obligations.
A future Fej\'er session can resume from the existing PAPER note in
Section~\ref{sec:fejer}; reconstruction, certified evaluation and the signed
connection to the ledger must remain distinct. The tail, geometry,
coefficient and connecting sources still require independent compilation,
followed by certification of the primitives and quadratures actually chosen.
No new Fej\'er research or Lean theorem is introduced by this documentary
revision. A signed bound on the ledger remains an open obligation.

\clearpage
\section{Numerical addendum at $N=10^8$}\label{app:numerical}
\textbf{Observed native arithmetic; 4 October 2026.} The finite run uses
\[
N=100\,000\,000,\quad K=134\,217\,728=2^{27},\quad
S=288\,230\,376\,151\,711\,744=2^{58}.
\]
The exact reconstructed integer is
\[
\boxed{I_N=14616354094337692277256326339388864358034646.}
\]
The five residues, in the modulus order of Section~\ref{sec:crt}, are
\[
(577099750,\ 626366149,\ 232952773,\ 2618309304,\ 2671748987).
\]
An independent Python CRT reconstruction gives the same integer as the
producer, the native direct $A32$ sum and the auxiliary $B40$ reconstruction:
\[
I_{\mathrm{CRT}}=I_{\mathrm{direct},A32}=I_{B40}=I_N.
\]
Its normalized value is
\[
I_N/S^2=175937962.675324617193216257686209932693\ldots.
\]
This number is the normalized canonical integer, not an exact evaluation of
the real logarithmic coefficient $C_N$ or of a truncated Mellin integral.
The fixed primary radius from the canonical logarithm contract is
\[
E_{\log}=\frac{1844674425817699235409551617}
{83076749736557242056487941267521536}
\approx2.2204460714547736\times10^{-8}.
\]
The observed normalized difference is zero, and exact rational comparison
confirms
\[
0=|I_{A32}-I_{B40}|/S^2\le2E_{\log}
\approx4.4408921429095471\times10^{-8}<10^{-6}.
\]
The $C_N$ enclosure inherits the canonical correspondence assumptions;
the auxiliary $B40$ real-log refinement remains unproved in Lean.

\noindent\begin{tabularx}{\textwidth}{@{}Yrrr@{}}
\toprule
Process & Wall (s) & CPU (s) & Peak RSS (bytes)\\
\midrule
Producer BUILD04 & 1098.219 & 1097.65625 & 1743966208\\
Operational checker & 1659.281 & 1657.84375 & 1744252928\\
\bottomrule
\end{tabularx}
The corresponding parent wall times are 1099.156 and 1661.266 seconds.
The overlapping executions and finalization take 1678.252969 seconds
(27 minutes 58.253 seconds). Both receipts confirm driver exit, an empty
Job and validated POST observations. The peak RSS is about 1.625 GiB per
native process; it is distinct from the requested 3 GiB commit quota.

\textbf{Unchanged boundary.} There are still 88 credited Lean modules and
1\,488 auxiliary declarations, including definitions. Native Lean refinement,
the complete Mellin coefficient truncation chain, spectral $H1$, control of
$D_N$ and Goldbach remain unproved. No new theory was undertaken for this run.

\clearpage
\subsection*{New numerical provenance}
\begingroup\small
The following paths are relative to \path{Goldbach_Numerical_20261004}.
Full SHA-256 values, file sizes, exact fractions and execution metadata are
retained in \path{numerical_verdict.json} and \path{evidence_v3_1.json}.
The verdict string of the native checker is
\path{EXACT_INTEGER_PROJECTION_CHECKED_PENDING_PRIMITIVE_PROOF}.
It is not a proof-success flag for the open analytic obligations.

\begingroup\small\setlength{\tabcolsep}{3pt}
\noindent\begin{tabularx}{\textwidth}{@{}p{27mm}Yp{27mm}@{}}
\toprule
Item & Relative path & SHA prefix\\
\midrule
Producer receipt & \path{producer_run02/receipt.json} & \texttt{a9dfea2ffde7}\\
Checker receipt & \path{checker_run01/receipt.json} & \texttt{5f435444318f}\\
Build receipt & \path{checker_build_receipt.json} & \texttt{b647aa25994f}\\
Checker source & \path{checker_dif22_operational.cpp} & \texttt{589f61d71717}\\
Checker image & \path{checker_dif22_operational.exe} & \texttt{268c4a481691}\\
Factor payload & \path{payload_run02/factors.bin} & \texttt{32450e19fb3d}\\
Prime records & \path{payload_run02/records.bin} & \texttt{bf6bfeb2d843}\\
Producer claim & \path{payload_run02/producer.txt} & \texttt{aeee44e0e6cd}\\
\bottomrule
\end{tabularx}
\endgroup

The new numerical archive preserves the Windows backend, checker, scripts,
contracts, logs and receipts, including the initial infrastructure failure.
The unchanged producer image remains in the historical V3 corpus, with its
hash bound by the execution receipt.

The repository \href{https://github.com/seisner1967-hash/goldbach-horizon}
{\texttt{goldbach-horizon}} separates the original V3 publication from the
new delivery and numerical archive:\par\noindent
\path{publications/2026-10-goldbach-synthesis-v3-1}\par\noindent
\path{research/goldbach-numerical-20261004}\par
The manifest and restore helper recover all original bytes and SHA-256
values, including losslessly fragmented large payloads. No file is omitted.

The entry point \path{HANDOFF_FEJER_V3_1.txt} accompanies this synthesis and
its registry. Fej\'er retains PAPER status; this addendum closes the native
run without starting that future research session.
\endgroup

\subsection*{Local documentary sources}
The following local sources preserve the written arguments and their
limitations. No new external references or secondary citations were verified.
\begin{thebibliography}{9}\small
\setlength{\itemsep}{0pt}\setlength{\parskip}{0pt}
\bibitem{v2} Durand Serge, \emph{Finite Identities and Explicit Error Budgets
in a Framework for Binary Goldbach}, V2 synthesis, 1 October 2026.
Title, author and continuity recorded from the retained public source.
\bibitem{mono} Durand Serge, final monograph and local extraction
\path{monographie.txt}; conventions updated by the source framework
\path{Goldbach_Final_20261001/goldbach_synthesis.tex}.
\bibitem{ntt} ROLE4, \emph{Finite projection contract},
\path{round22/role4/circle_ntt_paper01/projection_contract22_final.md}.
\bibitem{cost} ROLE4, \emph{Mellin--Lambda feasibility},
\path{round22/role4/mellin_lambda_feasibility_paper01/feasibility22.md}.
\bibitem{gap} ROLE4, \emph{Audit of the spectral connection to $D_N$},
\path{round22/role4/continuous_dn_gap_paper01/dn_gap_audit22.md}.
\bibitem{weil} Independent Judge, \emph{Signed explicit-formula review},
\path{round22/judge5/signed_weil_source_review01/review22.md}.
\bibitem{orbit} ROLE4, \emph{Zero-orbit review},
\path{round22/role4/zero_orbit_source_review01/review22.md}.
\bibitem{fejer} ROLE1, \emph{True positivity of $D$ and reconstruction
of the retained blocks},
\path{round22/role1/spectral_next_paper01/spectral_next22.md}.
\end{thebibliography}
\end{document}
